\documentclass[11pt]{article}
\usepackage[a4paper,margin=25mm]{geometry}
\usepackage{amsmath,amssymb,amsthm,booktabs,array}
\usepackage[T1]{fontenc}
\usepackage[colorlinks=true,linkcolor=blue,urlcolor=blue,citecolor=blue]{hyperref}
\newtheorem{theorem}{Theorem}
\newtheorem{lemma}[theorem]{Lemma}
\newtheorem{corollary}[theorem]{Corollary}
\newcommand{\R}{\mathbb{R}}
\newcommand{\C}{\mathbb{C}}
\newcommand{\ii}{\iota_z}
\newcommand{\leanfile}{\href{https://github.com/the-omega-institute/trureturing/blob/c55f6f7c7d7e3cf46b0a7da9efbba687448b29aa/D5/S3/Geometry/ProbabilityTransport/WeightedProjectionMinors.lean}{\texttt{WeightedProjectionMinors.lean}}}
\title{Proof of the weighted-minor inequality in\\MathOverflow 469781}
\author{Haobo Ma\\\small ChronoAI Pte Ltd, Singapore
\and Wenlin Zhang\\\small National University of Singapore\\\small ChronoAI Pte Ltd, Singapore}
\date{3 November 2026}
\begin{document}
\maketitle
\begin{abstract}
We prove the displayed inequality in MathOverflow question 469781 \cite{mo},
which compares weighted squared
two-by-two minors of an orthonormal pair with squared three-by-three minors,
weighted by the largest edge weight of each triple. We prove the inequality
by contraction in exterior algebra and a finite threshold decomposition.
The proof gives both minimum and maximum bounds, permits signed weights and
complex vectors, and extends to arbitrary exterior degrees and several
orthogonal lifts. The orthonormal comparison follows from Lyons (2003);
Miller and Bistro\'n (2025) give a direct two-sided proof and settle the
linked distance-triangle conjecture.
\end{abstract}

\section{The question and its scope}
MathOverflow question 469781, posted by \texttt{meler} on 23 April 2024
\cite{mo}, asks whether the inequality in Theorem~\ref{thm:main} holds for
an orthonormal triple and a symmetric matrix with strictly positive entries.
It does. The same question was posted on Mathematics Stack Exchange on
19 April 2024 \cite{mse}. Only off-diagonal entries occur, so we write
$a_{ij}=a_{ji}$ for $i\ne j$ and do not impose a diagonal convention.

For $x,y,z\in\R^n$, put
\[
 p_{ij}=x_i y_j-x_j y_i,\qquad
 q_{ijk}=\det\begin{pmatrix}x_i&x_j&x_k\\y_i&y_j&y_k\\z_i&z_j&z_k\end{pmatrix}.
\]
\begin{theorem}\label{thm:main}
Let $n\ge3$, $\|z\|=1$, and $\langle z,x\rangle=\langle z,y\rangle=0$.
For arbitrary real symmetric off-diagonal weights $a_{ij}$,
\begin{equation}\label{eq:main}
 \sum_{i<j}a_{ij}p_{ij}^{\,2}
 \ \le\ \sum_{i<j<k}\max\{a_{ij},a_{ik},a_{jk}\}\,q_{ijk}^{\,2}.
\end{equation}
In particular, this proves the inequality for every real orthonormal triple
and every positive-entry symmetric matrix.
\end{theorem}

For the orthonormal case, this is a specialization of Lyons's
Theorem~6.2 \cite{lyons}: the projection determinantal measures of
$\operatorname{span}(x,y)$ and $\operatorname{span}(x,y,z)$ are stochastically
ordered. Their probabilities on pairs and triples are $p_{ij}^2$ and
$q_{ijk}^2$, respectively. Applying that order to the increasing event
of containing a pair from $F$ gives the family comparison below.
The argument below gives the needed comparison
directly, including its threshold step and weaker vector hypotheses.

A later direct proof is Miller and Bistro\'n \cite{mb1}, Lemma~10,
equation~(27), first posted on 18 September 2025: it establishes both
the minimum and maximum comparisons for complex orthonormal triples
and nonnegative weights. The corresponding locators in their
4 October 2026 revision \cite{mb4} are Lemma~13 and equation~(34).
Theorem~2 of \cite{mb1} also proves the distance-triangle conjecture
mentioned in the posing thread, as discussed below. Our purpose is to give
the short contraction argument and its exterior-degree extensions.

\section{Contraction and pair families}
Equip the exterior powers of $\R^n$ with the inner product for which
$e_{i_1}\wedge\cdots\wedge e_{i_r}$, $i_1<\cdots<i_r$, are orthonormal.
Write $L_z v=z\wedge v$ and let $\ii=L_z^*$ denote contraction.
Choose an orthonormal basis with first vector $z$. Every exterior form $b$
has a unique orthogonal decomposition
\[
 b=z\wedge c+d,\qquad c,d\in\Lambda(z^\perp).
\]
Here $\ii b=c$, so $\|\ii b\|\le\|b\|$.
For $u=x\wedge y$ and $b=z\wedge u$, the hypotheses of
Theorem~\ref{thm:main} give
\begin{equation}\label{eq:recover}
 \ii b=u,\qquad
 \|b\|^2=\|u\|^2=\sum_{i<j}p_{ij}^2=\sum_{i<j<k}q_{ijk}^2.
\end{equation}
Indeed, $u\in\Lambda^2(z^\perp)$, and $z\wedge u=x\wedge y\wedge z$.
In coordinates the first identity is
$\sum_k z_k q_{kij}=p_{ij}$, with $q$ extended as an alternating array.

For a family $F$ of unordered pairs from $[n]=\{1,\ldots,n\}$, define
\[
 N(F)=\{T\subseteq[n]: |T|=3\text{ and some }S\in F\text{ satisfies }S\subset T\}.
\]
Let $P_F$ and $P_{N(F)}$ be the orthogonal coordinate projections onto
these pairs and triples.

\begin{lemma}\label{lem:family}
Under the vector hypotheses of Theorem~\ref{thm:main}, every pair family satisfies
\begin{equation}\label{eq:family}
 \sum_{\{i,j\}\in F,\ i<j}p_{ij}^2
 \ \le\ \sum_{\{i,j,k\}\in N(F),\ i<j<k}q_{ijk}^2.
\end{equation}
\end{lemma}
\begin{proof}
The coefficient of $e_i\wedge e_j$ in $\ii b$ uses only coefficients
$q_{kij}$. If $\{i,j\}\in F$, every triple contributing to that coefficient
belongs to $N(F)$. Consequently
\[
 P_F\ii b=P_F\ii P_{N(F)}b.
\]
Both an orthogonal projection and contraction by a unit vector have norm
at most one. Thus, using \eqref{eq:recover},
\[
 \|P_Fu\|=\|P_F\ii P_{N(F)}b\|
 \le\|\ii P_{N(F)}b\|\le\|P_{N(F)}b\|.
\]
Squaring gives \eqref{eq:family}. No summation over ordered tuples is used:
each pair and each triple occurs once.
\end{proof}

\section{From families to weights}
\begin{proof}[Proof of Theorem~\ref{thm:main}]
First suppose $a_{ij}\ge0$. List their distinct positive values as
$0<t_1<\cdots<t_s$, and set $t_0=0$. For each level put
$F_\ell=\{\{i,j\}:a_{ij}\ge t_\ell\}$. Then
\[
 a_{ij}=\sum_{\ell=1}^s(t_\ell-t_{\ell-1})
                 \mathbf1_{\{i,j\}\in F_\ell}.
\]
Similarly, for $T=\{i,j,k\}$,
\[
 \max\{a_{ij},a_{ik},a_{jk}\}
 =\sum_{\ell=1}^s(t_\ell-t_{\ell-1})\mathbf1_{T\in N(F_\ell)}.
\]
Multiply Lemma~\ref{lem:family} for $F_\ell$ by
$t_\ell-t_{\ell-1}\ge0$ and sum over $\ell$. These two identities give
\eqref{eq:main}. If all weights vanish, both sides vanish and the level
sum is empty.

For general real weights, let $m=\min_{i<j}a_{ij}$ and apply the result to
$a'_{ij}=a_{ij}-m\ge0$. The triple maximum also decreases by $m$.
The two sides of \eqref{eq:main} therefore each decrease by the same
quantity $m\|u\|^2$, by \eqref{eq:recover}. Cancelling it proves the assertion.
\end{proof}

\begin{corollary}\label{cor:complex}
Theorem~\ref{thm:main} holds for $x,y,z\in\C^n$, with Hermitian
orthogonality and squared absolute values of the minors.
\end{corollary}
\begin{proof}
Use the Hermitian inner product on exterior powers. The orthogonal
decomposition $b=z\wedge c+d$ and the contraction bound are unchanged;
the coordinate contraction is now $\sum_k\overline{z_k}q_{kij}$.
Projection locality is identical. Lemma~\ref{lem:family}, the level sums,
and the constant shift consequently hold with absolute squares.
\end{proof}

\begin{corollary}[Two-sided comparison]\label{cor:twosided}
Under the hypotheses of Theorem~\ref{thm:main}, or its complex version,
\[
 \sum_{i<j<k}\min\{a_{ij},a_{ik},a_{jk}\}|q_{ijk}|^2
 \ \le\ \sum_{i<j}a_{ij}|p_{ij}|^2
 \ \le\ \sum_{i<j<k}\max\{a_{ij},a_{ik},a_{jk}\}|q_{ijk}|^2.
\]
\end{corollary}
\begin{proof}
Apply the upper bound to $-a_{ij}$ and use
$\max\{-a_{ij},-a_{ik},-a_{jk}\}=-\min\{a_{ij},a_{ik},a_{jk}\}$.
This recovers the two-sided inequality of \cite[Lemma~10]{mb1}
and extends it to signed weights and the weaker vector hypotheses.
\end{proof}

\section{All exterior degrees and repeated lifts}
Write $\binom{[n]}k$ for the $k$-element subsets of $[n]$; coefficients
of exterior forms always use increasing indices.
\begin{corollary}\label{cor:degrees}
Let $\mathbb K=\R$ or $\C$, let $1\le k$, $1\le r$, and $k+r\le n$.
Let $z_1,\ldots,z_r$ be orthonormal, set
$V=\operatorname{span}(z_1,\ldots,z_r)^\perp$, and take any
$u\in\Lambda^k V$, not necessarily decomposable. Put
$b=z_r\wedge\cdots\wedge z_1\wedge u$.
For every real array $(a_S)_{S\in\binom{[n]}k}$,
\begin{equation}\label{eq:degrees}
 \sum_{T\in\binom{[n]}{k+r}}\left(\min_{\substack{S\subset T\\|S|=k}}a_S\right)|b_T|^2
 \ \le\ \sum_{S\in\binom{[n]}k}a_S|u_S|^2
 \ \le\ \sum_{T\in\binom{[n]}{k+r}}\left(\max_{\substack{S\subset T\\|S|=k}}a_S\right)|b_T|^2.
\end{equation}
\end{corollary}
\begin{proof}
First take $r=1$, $z=z_1$. The decomposition used above works in every
degree: $\iota_z(z\wedge u)=u$ and $\|z\wedge u\|=\|u\|$ for
$u\in\Lambda^k(z^\perp)$. For any family $F\subseteq\binom{[n]}k$,
let $N(F)$ consist of the $(k+1)$-sets containing a member of $F$.
A coefficient of $\iota_z b$ at $S$ uses only coefficients $b_T$ with
$S\subset T$, so the same projection identity gives
$\sum_{S\in F}|u_S|^2\le\sum_{T\in N(F)}|b_T|^2$.
The threshold argument and equal-mass shift prove the upper bound for
real weights, and negation proves the lower bound.

For $r>1$ apply this one-step result successively while adjoining
$z_1,\ldots,z_r$. Each intermediate form is perpendicular in every
factor to the next vector, by orthogonality. Iterating the maximum
over subsets gives precisely $\max_{S\subset T,\,|S|=k}a_S$; iterating
the minimum gives the analogous minimum. This proves \eqref{eq:degrees}.
\end{proof}

\begin{corollary}[Constant weights and the three-dimensional case]\label{cor:sharp}
Constant weights give equality in \eqref{eq:degrees}.
The coefficient one in its nonnegative-weight upper bound is sharp.
For an orthonormal triple in dimension three, the original inequality
is the elementary bound in the posing questions.
\end{corollary}
\begin{proof}
Orthogonal lifts preserve the norm, so both total masses are equal.
Taking $u=e_1\wedge\cdots\wedge e_k$, $z_j=e_{k+j}$ and constant
positive weights gives a nonzero equality example. When $n=3$ and
$x,y,z$ are orthonormal, $|q_{123}|^2=\sum_{i<j}|p_{ij}|^2=1$,
and bounding each weight by the single triple maximum is the proof
already given in \cite{mo,mse}.
\end{proof}

\section{Coverage and verification}
The comment thread on Mathematics Stack Exchange relates the minors
question to the triangle inequality for weighted exterior-product distances
\cite{mse}. That larger question is formulated as Conjecture~2.1 in the
2023 version of Bistro\'n et al.\ \cite{distances}. It has a later full proof:
Miller and Bistro\'n \cite[Theorem~2]{mb1} prove that, for every finite
distance matrix $(d_{ij})$ and $p\ge2$,
\[
 D_p(x,y)=\left(\sum_{i<j}d_{ij}^{\,p}|x_i y_j-x_j y_i|^2\right)^{1/p}
\]
is a distance on complex projective space, with unit representatives $x,y$.
The conjecture is the $p=2$ instance. This is Theorem~1 in \cite{mb4}.
The complete distance theorem is cited here, while the written proofs above
establish the exterior comparisons. In the $p=2$ specialization their
weights are $a_{ij}=d_{ij}^2$.
Table~\ref{tab:coverage} records the source items and their mathematical scope.

\begin{corollary}[Sharp exponent for the linked distance]\label{cor:exponent}
For every $n\ge2$ and every distance matrix, $D_p$ fails the triangle
inequality when $0<p<2$. For constant off-diagonal distances,
$D_p$ is a distance for every $p\ge2$.
\end{corollary}
\begin{proof}
Take $x=\cos\theta\,e_1+\sin\theta\,e_2$,
$y=\cos\theta\,e_1-\sin\theta\,e_2$, and $z=e_1$.
For $0<\theta<\pi/2$, their three distances are
$d_{12}\sin^{2/p}(2\theta)$ and twice $d_{12}\sin^{2/p}\theta$.
The triangle inequality would require
$\cos\theta\le2^{p/2-1}$, which fails for sufficiently small positive
$\theta$ if $p<2$. This is the obstruction in \cite[Remark~4]{mb1,mb4}.
If all $d_{ij}=c>0$, the Gram identity gives
$D_p(x,y)=c(\|xx^*-yy^*\|_{\mathrm{HS}}/\sqrt2)^{2/p}$.
The Hilbert--Schmidt norm obeys the triangle inequality, and
$(a+b)^s\le a^s+b^s$ for $a,b\ge0$ and $0<s\le1$, proving the claim.
\end{proof}

\begin{table}[ht]
\centering\small
\begin{tabular}{@{}p{0.43\linewidth}p{0.52\linewidth}@{}}
\toprule
Source item & Coverage \\
\midrule
MathOverflow 469781, displayed question & Proved: Theorem~\ref{thm:main}. Prior proof cited: \cite[Theorem~6.2]{lyons}, then \cite[Lemma~10]{mb1}.\\
Both posing posts, $n=3$ calculation & Prior proof cited: the questions themselves; Corollary~\ref{cor:sharp}.\\
Mathematics Stack Exchange 4902019, displayed question & Proved: the identical inequality, Theorem~\ref{thm:main}.\\
Same thread, 22 April 2024 comments; \cite{distances}, Conjecture~2.1 & Prior proof cited: \cite[Theorem~2]{mb1}, with $p=2$; \cite[Theorem~1]{mb4}.\\
\cite{distances}, Corollary~4.2 and Theorems~6.4, 6.8 & Prior proofs cited: dimension three, Euclidean and multiplicative distances, respectively. All are included in \cite[Theorem~2]{mb1}.\\
\cite{mb1}, Lemma~10, equation~(27); \cite{mb4}, Lemma~13, equation~(34) & Prior proof cited; Corollary~\ref{cor:twosided} supplies the signed and weaker-hypothesis version.\\
\cite{mb1,mb4}, Remark~4; constant distances & Prior proof cited and proved: sharpness of $p\ge2$ and the constant-distance case, Corollary~\ref{cor:exponent}.\\
Signed, complex, higher-degree and repeated-lift comparisons & Proved: Corollaries~\ref{cor:complex}--\ref{cor:degrees}.\\
\bottomrule
\end{tabular}
\caption{The posing posts and their linked mathematical question.}
\label{tab:coverage}
\end{table}

The declaration \texttt{weighted\_minors} in the namespace
\texttt{AlternatingArrayContraction} of the pinned file \leanfile\ states
the literal real orthonormal-triple inequality with strictly positive
symmetric entries, namely that instance of Theorem~\ref{thm:main}.
Its determinant definitions and increasing-index sums agree with the ones
above. The signed-weight, weaker-hypothesis, complex, higher-degree and
repeated-lift extensions have the written proofs above, rather than that
declaration's narrower scope. The accompanying \texttt{verify.py} supplies
exact rational checks of pair families and signed weights in dimensions
three through six, and exact Gaussian-rational checks of contraction,
two-sided bounds, higher degrees and repeated lifts in dimensions four
through six. These finite comparisons are distinct from the universal proofs.

\begin{thebibliography}{9}
\small
\bibitem{mo}
\texttt{meler}, \emph{Inequality involving minors of an orthogonal matrix},
MathOverflow question 469781, 23 April 2024.
\url{https://mathoverflow.net/q/469781}.

\bibitem{mse}
\texttt{meler}, \emph{Inequality involving a symmetric matrix and minors of
an orthogonal matrix}, Mathematics Stack Exchange question 4902019,
19 April 2024; comments dated 22 April 2024.
\url{https://math.stackexchange.com/q/4902019}.

\bibitem{lyons}
R. Lyons, \emph{Determinantal probability measures},
Publications Math\'ematiques de l'IH\'ES \textbf{98} (2003), 167--212,
especially Theorem~6.2, p.~183.
\href{https://doi.org/10.1007/s10240-003-0016-0}{doi:10.1007/s10240-003-0016-0}.

\bibitem{distances}
R. Bistro\'n, M. Eckstein, S. Friedland, T. Miller, and K. \.Zyczkowski,
\emph{A new class of distances on complex projective spaces},
\href{https://arxiv.org/abs/2312.02583v1}{arXiv:2312.02583v1} (2023),
Conjecture~2.1. Published in Linear Algebra and its Applications
\textbf{721} (2025), 577--611,
\href{https://doi.org/10.1016/j.laa.2024.10.017}{doi:10.1016/j.laa.2024.10.017}.

\bibitem{mb1}
T. Miller and R. Bistro\'n,
\emph{Distances between pure quantum states induced by a distance matrix},
\href{https://arxiv.org/abs/2509.14727v1}{arXiv:2509.14727v1},
18 September 2025, Theorem~2, Theorem~9 and Lemma~10, equation~(27).

\bibitem{mb4}
T. Miller and R. Bistro\'n,
\emph{Distances between pure quantum states induced by a distance matrix},
\href{https://arxiv.org/abs/2509.14727v4}{arXiv:2509.14727v4},
4 October 2026, Theorem~1, Theorem~12 and Lemma~13, equation~(34).
\end{thebibliography}

\bigskip\noindent\small
\textcopyright\ 2026 Haobo Ma and Wenlin Zhang.
Licensed under \href{https://creativecommons.org/licenses/by/4.0/}{Creative Commons Attribution 4.0 International (CC BY 4.0)}.
\end{document}
